% ============================================================
\section{Where Responsibility Lands}\label{sec:responsibility}
% ============================================================

The Cascade identifies which layer failed and how the failure propagated; the downstream question of who bears legal responsibility requires a doctrinal methodology (developed in companion work). The Cascade does make visible a structural pressure on governance institutions that is worth stating independently.

\subsection{The Upstream Pressure}

Output-only review cannot stabilize same-case-same-action, because the reviewer at the epistemic boundary sees the output trace, not the internal states that produced it (\S\ref{sec:cascade}). The Cascade does not resolve the downstream legal question---that requires a doctrinal methodology developed in companion work. But it makes visible a structural pressure that legal analysis must address. When output-only review cannot stabilize same-case-same-action, responsibility drifts upstream---the responsibility gap that the ethics literature names in the abstract~\cite{konigs2022} and that clinical decision-support studies observe as a diffusion of accountability across operator, deployer, and developer~\cite{bleher2022}---toward entities whose decisions are institutionally legible (product design, organizational resourcing, audit obligations) rather than entities whose fault is partitioned from observation by the epistemic boundary. The drift has an addressee problem---in cases (a) and (b), the victims cannot bear liability (a minor lacks legal capacity; a person of diminished capacity lacks the agency any liability regime presupposes)---but whether the resulting upstream allocation is doctrinally sound is a legal question, not a taxonomic one. The Cascade provides the diagnostic vocabulary that legal analysis requires; it does not substitute for it.

\subsection{Institutional Implications}

The Cascade maps each layer to a distinct institutional response domain (Table~\ref{tab:interventions}). The point is not that law redeems technical failure. It is that technical systems cannot settle their own authority boundaries once the structural insufficiency at the epistemic boundary is in view. The Cascade provides the diagnostic vocabulary for institutions to ask, and answer, \textit{which layer failed}---separating that architectural question from the downstream allocation of responsibility.

\begin{table}[!htbp]
\centering
\small
\caption{Layer-specific institutional interventions.}
\label{tab:interventions}
{\setlength{\tabcolsep}{4pt}
\begin{tabularx}{\columnwidth}{@{}P{1.8cm}P{2.7cm}Y@{}}
\toprule
\textbf{Layer} & \textbf{Intervention domain} & \textbf{Mechanism} \\
\midrule
L1 Environment & Market-level & Labeling of intended user context; interoperability mandates; data-provenance documentation for training corpora \\
L2 Interface & Session architecture & Context transparency; session boundary instrumentation; crisis-detection mandates for extended interactions \\
\textit{Epistemic Gap} & \textit{Audit infrastructure} & \textit{Logging mandates; structural-access requirements for certified auditors; min-cut computation on deployment topology} \\
L3 Algorithm \& Weights & Training governance & Weight-level audit protocols; RLHF methodology transparency; red-team requirements; sycophancy and contamination diagnostics \\
L4 Strategic Authority & Strategy governance & Longitudinal re-certification; cross-version behavioral audits; mandatory disclosure of safety-resource allocation \\
\bottomrule
\end{tabularx}}
\end{table}

